% frontmatter/abstract.tex -- ABSTRACT (thesis.md S3.11): starts with
% the English translation of the thesis title, then the word ABSTRACT.

\starChapterHead{{\bfseries\itshape\MakeUppercase{\thesisTitleEnglish}\par}\vskip 1.5\baselineskip}{ABSTRACT}
\addcontentsline{toc}{chapter}{ABSTRACT}

\noindent Indonesian-language product reviews on e-commerce platforms keep
growing every day, making manual assessment of consumer opinions
inefficient. This research aims to design, build, and evaluate a
sentiment classification model based on Long Short-Term Memory
(LSTM) to classify product reviews into three classes, namely
positive, negative, and neutral, and to compare its performance
against a conventional machine learning method based on Support
Vector Machine (SVM) as a baseline. The method includes data
preprocessing tailored to the morphological characteristics of the
Indonesian language, such as informal-word normalization and
stopword removal, followed by feature extraction using word
embeddings and training an LSTM model with a single hidden layer of
128 units. Model performance is evaluated using accuracy, precision,
recall, and F1-score on a held-out test set. The test results show
that the proposed LSTM model achieves an F1-score of 0.85, higher
than the baseline SVM model, which only achieves an F1-score of
0.75. The performance improvement is particularly evident in the
LSTM model's ability to handle sentences containing negation,
although the model still shows limitations in handling reviews that
contain sarcasm or implicit opinions. The results of this research
are expected to serve as a basis for developing an automated
consumer-opinion monitoring system.

\vskip 2\baselineskip
\noindent\textbf{Keywords : sentiment analysis, LSTM, product reviews,
Support Vector Machine}
